\section{Analysen}
\label{sec:Analysen}

\subsection{Ist-Analyse}
\label{sec:Ist-Analyse}
%in den verschiedenen Status (Planung, Antrag, genehmigt), welche ebenfalls halber Urlaubstage und Sonderurlaub beinhalten 
Derzeit verwaltet unser Backoffice für jeden Mitarbeiter alle Krankheitstage, Urlaubsanträge sowie Schulblöcke der Auszubildenden in einem tagesgenauen Excel-Dokument. Die Urlaubsanträge, welche ebenfalls halbe Urlaubstage und Sonderurlaub beinhalten, werden mit den verschiedenen Status (Planung, Antrag, genehmigt) eingetragen.
Das Dokument wird verwendet, um Schulungen oder Projekte zu planen, Urlaubszeiträume abzustimmen und um gegenüber dem Vorstand stets aussagefähig zu sein, wenn dieser Fragen zu Abwesenheiten hat. 

Für die Verwaltung interner Angelegenheiten wird das MIBS virtualDesk verwendet. Die Plattform besteht aus einer SAPUI5/Fiori Elements\footnote{\label{Fiori}SAP Framework für UI5-basierte Anwendungen} Web-Oberfläche und einem SAP S/4HANA\footnote{ERP-Suite basierend auf der In-Memory-Datenbank HANA} Backend. Dieses verfügt über Apps, mit welchen unter anderem Urlaubsanträge und Krankmeldungen gepflegt werden können. Auch die Zeiträume der Schulblöcke für die Auszubildenden werden bereits von der Ausbildungsleitung im MIBS virtualDesk gepflegt. Die Informationen aus diesen Apps werden genutzt, um die Excel-Datei täglich zu aktualisieren. Dazu müssen die verschiedenen Apps des MIBS virtualDesk aufgerufen und die Informationen manuell in das Dokument übertragen werden. Gleichzeitig wird die Tabelle neben dem MIBS virtualDesk geöffnet, um während der Bearbeitung einen Überblick über mögliche Terminüberschneidungen zu erhalten. Die Liste wird gleichzeitig von der Ausbildungsleitung genutzt, um anstehende Ausbildungsinhalte zu planen. Die tägliche Pflege und der Abgleich beansprucht wertvolle Zeit des Backoffice-Teams.
% beansprucht oder beanspruchen?
% TODO Vorschlag überprüfen
% Die MIBS AG verwendet für ihre internen, administrativen und mitarbeiterbezogenen Prozesse eine browserbasierte Plattform, das sogenannte MIBS virtualdesk (kurz MVD). Das MVD besteht aus einer SAPUI5/Fiori Web-Oberfläche und einem SAP S/4HANA Backend. Es verfügt über Apps, in denen Mitarbeiter- und Kundendaten hinterlegt, Kundeneinsätze und -projekte dokumentiert sowie Arbeitszeiten und sämtliche An- und Abwesenheiten (Krankheit, Urlaub, Sonderurlaub, Berufsschule) erfasst werden. 
% Es gibt allerdings keine App, die die An- und Abwesenheitsdaten der Mitarbeiter übersichtlich zusammenfasst und anzeigt, so dass das Backoffice gezwungen ist, eine Excel-Tabelle zu führen, in die diese Daten täglich manuell aus dem MVD übertragen werden. Hierzu müssen nacheinander die jeweiligen Apps (Urlaub, Krankheit) geöffnet und die tagesaktuellen Daten pro Mitarbeiter in die Exceltabelle übertragen werden. Die tägliche Pflege und der Abgleich beansprucht wertvolle Zeit des Backoffice-Teams.


\subsection{Wirtschaftlichkeitsanalyse}
\label{sec:Wirtschaftlichkeitsanalyse}

\subsubsection{\gqq{Make-or-Buy} Entscheidung}
\label{sec:Make-or-Buy Entscheidung}
Bei der im Betrieb verwendeten Software, dem MIBS virutalDesk, handelt es sich um eine eigens entwickelte Individualsoftware. Der Markt bietet hierfür keine passende Lösung, die alle Anforderungen erfüllt. Aus diesem Grund soll das Projekt als interne Entwicklung realisiert werden.

\subsubsection{Projektkosten}
\label{sec:Projektkosten}

Im Folgenden wurden die voraussichtlich anfallenden Kosten kalkuliert. Die Projektkosten setzen sich aus den Personalkosten und dem Gemeinkostenzuschlag zusammen. 
Die Personalkosten sind folglich aufgeteilt: Auf den Auszubildenden, welcher für die Umsetzung des Projektes zuständig ist, den Projektbetreuer, der das Codereview durchführt und das Projekt beratend betreut sowie die Bürofachkraft des Backoffice, welche in die Anwendung eingearbeitet werden muss.
Bei den Gehältern handelt es sich um Mittelwerte, die von der Personalabteilung bereitgestellt worden sind. Über den Gemeinkostenzuschlag werden alle Kosten erfasst, die nicht spezifisch zuzuordnen sind. Unter diese fallen beispielsweise Strom- oder Verwaltungskosten.
%TODO fußnote Fakura - sowas wie: Arbeitszeit, die unmittelbar an einem Kundenprojekt aufgewendet wird

\tabelle{Kostenrechnung}{tab:Kostenrechnung}{Kostenrechnung}
\newpage

\subsubsection{Amortisationsdauer}
\label{sec:Amortisationsdauer}

Bevor die Amortisationsdauer berechnet werden kann, muss zunächst der aktuelle Zeitaufwand ermittelt werden. Hierfür wurde über die Dauer einer Arbeitswoche die Höhe des Zeitaufwands dokumentiert. In der Erfassung wurde ausschließlich die Pflege der Daten berücksichtigt, da dieser Arbeitsschritt komplett entfallen soll. Wie aus \TabelleRef{tab:Wöchentlicher Pflegeaufwand} zu entnehmen ist, beträgt der durchschnittliche, wöchentliche Aufwand zum Pflegen der Daten 15 Minuten pro Tag.

\tabelle{Wöchentlicher Pflegeaufwand}{tab:Wöchentlicher Pflegeaufwand}{Wöchentlicher Pflegeaufwand} 
% Fußnote Gemeinkostenzuschlag

Für die Ermittlung der Amortisationsdauer, wird zunächst die jährliche Einsparung berechnet. Diese ergibt sich aus dem jährlichen Arbeitsaufwand sowie den Personalkosten. Der jährliche Aufwand definiert sich durch die Arbeitstage multipliziert mit dem täglichen Aufwand. Die jährlichen Arbeitstage belaufen sich nach Abzug der Betriebsferien auf 246 Tage\footnote{Arbeitstage in der MIBS AG für 2024}.

\[
0{,}25 \, \frac{\text{Stunde}}{\text{Tag}} \times 50 \, \frac{\text{€}}{\text{Stunde}} \times 246 \, \frac{\text{Tage}}{\text{Jahr}} = 3075 \, \frac{\text{€}}{\text{Jahr}}
\]

% TODO ggf. auf nächste Seite
Die Amortisationsdauer ergibt sich wie folgt:

\[
\text{Amortisationsdauer} = \frac{\text{Projektkosten}}{\text{Jährliche Einsparung}}
\]

\[
\frac{3245\,\text{€}}{3075\,\text{€/Jahr}} \approx 1{,}055 \, \text{Jahre}
\]
Wie in \Abbildung{Amortisationsdauer} ersichtlich, zeigt sich, dass sich das Projekt nach etwa 1,055 Jahren beziehungsweise rund 12,7 Monaten refinanziert.


\begin{figure}[htb]
\centering
\includegraphicsKeepAspectRatio{Amortisationsdauer.png}{0.7}
\caption{Amortisationsdauer}
\label{fig:Amortisationsdauer}
\end{figure}

\subsubsection{Nicht monetäre Vorteile}
\label{sec:Nicht monetäre Vorteile}
Neben den kalkulierbaren Kostenersparnissen ergeben sich durch dieses Projekt auch nicht monetarisierbare Vorteile. Durch das gezielte Wegfallen von redundanten Aufgaben, beispielsweise der täglichen Datenpflege, kann mit einer Steigerung der Arbeitsmoral gerechnet werden. 
Des Weiteren wird das Risiko der Übertragung fehlerhafter Daten ausgeschlossen. 
Nicht zuletzt trägt es zu der Digitalisierung, Modernisierung und Vereinheitlichung der internen Prozesse im Unternehmen über das MIBS virtualDesk bei.


\subsection{Risikoanalyse}
\label{sec:Risikoanalyse}
Vor Projektbeginn wurde eine Risikomatrix gemäß ISO 31010 erstellt, um die potenziellen Risiken zu bewerten. Die X-Achse zeigt die Schwere der Folgen, die Y-Achse die Eintrittswahrscheinlichkeit. Risiken im \textcolor{Green}{grünen} Bereich gelten als unkritisch, während \textcolor{Goldenrod}{gelbe} Risiken Anpassungen in der Planung erfordern. Risiken im \textcolor{red}{roten} Bereich stellen eine erhebliche Gefährdung des Projektplans dar und erfordern ein gezieltes Risikomanagement. Die Matrix dient als Grundlage für die Risikobewertung durch die Projektverantwortlichen.\\ 
\newpage
Als plausible Risiken werden erkannt:
\begin{enumerate}
\item \textbf{Krankheit des Entwicklers:\\}
Da die Entwicklung von einem einzelnen Entwickler durchgeführt wird, kann ein krankheitsbedingter Ausfall kritisch sein. Der geplante Aufwand enthält jedoch Puffer, um Verzögerungen durch intensivere Arbeit an anderen Projekttagen auszugleichen. Aufgrund vergangener Erfahrungen, wird dieses Risiko als gering eingestuft. Eine lückenlose Dokumentation aller Prozesse und Entwicklungsschritte stellt sicher, dass bei längerer Abwesenheit des Entwicklers eine reibungslose Übergabe möglich ist.
\item \textbf{Ausfall der Hardware:\\}
Ein Ausfall der Hardware wird, basierend auf Erfahrungswerten, als unwahrscheinlich eingestuft. Im Falle eines Ausfalls stehen Ersatzgeräte der MIBS AG zur Verfügung, um Verzögerungen zu vermeiden.
\item \textbf{SAP-Systemausfall:\\}
Der Ausfall des SAP-Systems hätte erhebliche Auswirkungen, da das Backend vollständig auf diesem System basiert. Dieses Risiko wird jedoch als unwahrscheinlich eingeschätzt. Teile des Frontends können unabhängig vom System lokal entwickelt werden, wodurch die Auswirkungen begrenzt werden.
\item \textbf{Ausfall des Fachbereiches:\\}
Ein vollständiger Ausfall des Fachbereichs wird aufgrund der personellen Besetzung als unwahrscheinlich betrachtet. Die Auswirkungen wären gering, da die Anforderungen bereits definiert und die fachlichen Prozesse durchdacht sind.
\end{enumerate}

\begin{figure}[htb]
\centering
\includegraphicsKeepAspectRatio{Risiko.png}{0.9}
\caption{Risikoanalyse}
\label{fig:Risikoanalyse}
\end{figure}

\subsection{Datenanalyse}
\label{sec:Datenanalyse}
Für die Datenanalyse wurde das bestehende Datenmodel untersucht. Dabei wurden die benötigten Views identifiziert.
Eine Beschreibung der Views sowie deren Zusammenspiel wird im \Anhang{app:Datenanalyse} genauer erläutert.

%TODO: Tag statt Periode erwähnen? Bzw Daten hier beschreiben